\chapter{Conclusion}\label{ch:conclusion}

The study of sub-Neptune exoplanets occupies a critical intersection between geophysics and
atmospheric science. This work has demonstrated that the release of hydrogen from a miscible
sub-Neptune interior cannot be timed from equilibrium miscibility data alone. Between the
equilibrium onset of exsolution and the loss of local stability lies a metastable window,
and its width is set by a part of the free energy that the existing simulations do not
sample. By utilizing the \Atm{} framework to compute binodals, spinodals and their
uncertainties for H$_2$--MgSiO$_3$ parcels across local hydrogen mass fractions of 0.1--20\%
and pressures of 2--10\,GPa, we have constrained the thermodynamic limits on hydrogen
exsolution during the cooling of the magma ocean. This chapter summarizes the primary
insights derived from this modeling and discusses their implications for interpreting
observations.

\section{Summary of Findings}

Our modeling yields four primary conclusions regarding the exsolution of hydrogen from
sub-Neptune interiors.

\subsection{A Composition-Consistent Mixing Model}

We implemented the final 2026 hydrogen--silicate mixing model in \Atm{}. The published model
uses oxide units, whereas \Atm{} uses MgSiO$_3$ formula units, and the conversion must act on
the extensive free energy, not only on the compositions. With this conversion our
implementation reproduces the published phase diagram with an RMS residual of \digRms{}\,K.
Without it,
the metastable interval of the baseline parcel is about nine times too short.

\subsection{The Metastable Window}

With the published fit, a parcel with 3\,wt\% local hydrogen at 2\,GPa enters the miscibility
gap at \pubTbTwo{}\,K and becomes unstable at \pubTspTwo{}\,K. Over this \pubIntTwo{}\,K
interval, delayed exsolution could withhold up to \pubFhTwo{}\% of its hydrogen from
separation. In the published fit, parcels with less than \pubThreshold{}\,wt\% hydrogen never
become unstable, so their separation always requires nucleation; in the refits this threshold
lies lower.

\subsection{The Source Data Do Not Fix the Window}

The published fit departs from its own source simulations in two places: the compositions of
the two runs closest to the critical curve, where the source table and figure also disagree,
and the excess volume, which is about half the value that the fitted pressure scale implies.
Refitting the same model form to all source data shortens the interval to
\IntMOneTable{}\,K (68\% range \IntMOneTableLo{}--\IntMOneTableHi{}\,K). Across four smooth
model classes and six data treatments, the median interval ranges from \IntRangeLo{} to
\IntRangeHi{}\,K, and the most flexible class still allows the published value. The window
depends on the asymmetry, the excess entropy and the strength of the interaction together, and
the coexisting compositions and the homogeneous simulations pull these in different
directions.

\subsection{The Spinodal Is Not Identifiable from Coexistence Data}

A free-energy perturbation that leaves every equilibrium observable unchanged still moves the
spinodal by tens to hundreds of kelvin. Any spinodal inferred from coexistence data is
therefore a property of the assumed functional form. The curvature inside the miscibility gap
has to be measured, not inferred.

\section{Implications for Observations}

These findings introduce a new uncertainty into the interpretation of young sub-Neptunes.
Evolution models with hydrogen--silicate miscibility time the release of hydrogen with the
equilibrium binodal. If nucleation is suppressed, the release can lag by up to several hundred
kelvin of cooling and then occur abruptly.

Therefore, observational constraints from instruments such as JWST should be interpreted
within a coupled framework:
\begin{enumerate}
\item \textbf{Radius evolution:} Inflation of the envelope by exsolved hydrogen may set in later
than equilibrium models predict. The demographic tests proposed for young sub-Neptunes
\cite{rogers2026} should therefore include the metastable window as a systematic
uncertainty.
\item \textbf{Atmospheric composition:} A delayed and abrupt supply of hydrogen would
temporarily lower the mean molecular weight and shift the carbon speciation toward CH$_4$, in
line with the hydrogen-budget dependence found in our thesis \cite{laga2025}. Such a signal
requires efficient transport from the interior and remains untested.
\item \textbf{Interpretive caution:} Because the window is not identifiable from present data,
atmospheric or radius signatures should not yet be interpreted as evidence for or against
spinodal decomposition.
\end{enumerate}

\section{Future Outlook}

While this work establishes a thermodynamic bound on delayed hydrogen exsolution,
the sub-Neptune environment remains complex. The current models assume a binary
H$_2$--MgSiO$_3$ mixture with temperature-independent excess enthalpy and entropy. Future
simulations should measure the curvature of the homogeneous free energy inside the miscibility
gap directly, through the concentration structure factor, and test the temperature dependence
of the excess properties below $T_C$. The interfacial tension and mobility then fix nucleation
rates, which an evolution model can use instead of the two limiting cases studied here.
Lastly, as the magma ocean solidifies, the exsolution of incompatible volatiles could lead to
late-stage atmospheric inflation. This, too, should be incorporated as a time-dependent
evolution of the cooling and crystallizing magma ocean, coupled to the atmospheric chemistry
of \Atm{}. Despite these open questions, this work highlights that accurately predicting when
a sub-Neptune releases its hydrogen requires a fundamental understanding of the metastable
states of its interior.
